% !TEX program = xelatex
% ================================================================
% 题型八——解析几何与直线关系
% 对应原卷：第 28 题      对应幻灯片：第 9 页
% 编译：xelatex 题目08_解析几何与直线关系.tex
% ================================================================
\documentclass[12pt,a4paper]{article}
\usepackage[margin=2.5cm]{geometry}
\usepackage{xeCJK}
\usepackage{amsmath,amssymb}
\usepackage{xcolor}
\setCJKmainfont{Microsoft YaHei}[BoldFont = Microsoft YaHei Bold]
\setlength{\parindent}{0pt}
\pagestyle{empty}
% 题型标题：深蓝色加粗（与 PPT 一致）
\newcommand{\qtitle}[1]{{\Large\bfseries\color[HTML]{1F3864}#1}\par\vspace{0.9em}}

\begin{document}

\qtitle{题型八——解析几何与直线关系}

若$a\in\mathbb{R}$，直线$l_1:x+2ay-1=0$，\par
$l_2:(3a-1)x-ay-1=0$，\par
则“$a=0$”是“$l_1\parallel l_2$”的（　）\par
\vspace{0.9em}
A．充分不必要条件 \qquad\qquad B．必要不充分条件\par
\vspace{0.45em}
C．充要条件 \qquad\qquad D．既不充分也不必要条件\par
\vspace{0.45em}

\end{document}
